\usepackage{amsmath,amssymb,amsthm}
\usepackage{xcolor}
\usepackage{mathtools}
\usepackage{enumitem}
\usepackage{bbm}
\usepackage{semantic}
\usepackage{listings}
\usepackage{courier}
\usepackage{tikz}
\usetikzlibrary{calc}
\usepackage{color}
\usepackage{fancyhdr}
\usepackage{lastpage}

\definecolor{codegreen}{rgb}{0,0.6,0}
\definecolor{codegray}{rgb}{0.5,0.5,0.5}
\definecolor{codepurple}{rgb}{0.58,0,0.82}
\definecolor{backcolour}{rgb}{0.95,0.95,0.92}
 
\lstdefinestyle{mystyle}{
    backgroundcolor=\color{backcolour},   
    commentstyle=\color{codegreen},
    keywordstyle=\color{magenta},
    numberstyle=\tiny\color{codegray},
    stringstyle=\color{codepurple},
    basicstyle=\footnotesize,
    breakatwhitespace=false,         
    breaklines=true,                 
    captionpos=b,                    
    keepspaces=true,                 
    numbers=left,                    
    numbersep=5pt,                  
    showspaces=false,                
    showstringspaces=false,
    showtabs=false,                  
    tabsize=2,
    basicstyle=\footnotesize\ttfamily
}

\lstset{style=mystyle}

\usepackage{hyperref}

\DeclareMathOperator*{\E}{\mathbb{E}}
\let\Pr\relax
\DeclareMathOperator*{\Pr}{\mathbb{P}}

\newcommand{\eps}{\epsilon}
\DeclareMathOperator*{\argmax}{arg\,max}
\DeclareMathOperator*{\argmin}{arg\,min}
\newcommand{\inprod}[1]{\left\langle #1 \right\rangle}
\newcommand{\R}{\mathbb{R}}

\newcommand{\norm}[1]{\left\lVert#1\right\rVert}
\newcommand{\abs}[1]{\left\lvert#1\right\rvert}
\newcommand{\paren}[1]{\left(#1\right)}
\newcommand{\sqb}[1]{\left[#1\right]}
\newcommand{\Mod}[1]{\ (\mathrm{mod}\ #1)}
\newcommand{\setmath}[1]{\left\{#1\right\}}
\newcommand*\xor{\mathbin{\oplus}}
\newcommand{\ceil}[1]{\lceil#1\rceil}
\newcommand{\floor}[1]{\lfloor#1\rfloor}
\newcommand{\bigoh}[1]{\mathcal{O}\paren{#1}}

\newenvironment{proofoutline}
 {\renewcommand\qedsymbol{}\proof[Proof outline]}
 {\endproof}

\newtheorem{theorem}{Theorem}
\newtheorem*{lemma*}{Lemma}
\newtheorem{corollary}[theorem]{Corollary}
\newtheorem{remark}{Remark}
% \newtheorem{lemma}[theorem]{Lemma}
\newtheorem{lemma}{Lemma}
% \newtheorem{example}[theorem]{Example}
\newtheorem{example}{Example}
\newtheorem{observation}[theorem]{Observation}
\newtheorem{proposition}[theorem]{Proposition}
% \newtheorem{definition}[theorem]{Definition}
\newtheorem{definition}{Definition}
\newtheorem{claim}[theorem]{Claim}
% \newtheorem{fact}[theorem]{Fact}
\newtheorem{fact}{Fact}
\newtheorem{assumption}[theorem]{Assumption}
\newtheorem{conjecture}[theorem]{Conjecture}
% \newtheorem{problem}[theorem]{Problem}
\newtheorem{problem}{Problem}
\newenvironment{solution}
  {\begin{proof}[Solution]}
  {\end{proof}}

\newcommand*\diff{\mathop{}\!\mathrm{d}}
\newcommand*\Diff[1]{\mathop{}\!\mathrm{d^#1}}

\DeclareMathOperator{\Tr}{Tr}
\DeclareMathOperator{\prox}{prox}
\DeclareMathOperator{\var}{var}
\DeclareMathOperator{\dom}{dom}
\DeclareMathOperator{\rk}{rk}
\DeclareMathOperator{\nul}{nul}
\DeclareMathOperator{\sign}{sign}
\DeclareMathOperator{\dual}{dual}

\newcommand{\exv}[1]{\E\sqb{#1}}
\newcommand{\prv}[1]{\Pr\sqb{#1}}

\newcommand*{\vertbar}{\rule[-1ex]{0.5pt}{2.5ex}}
\newcommand*{\horzbar}{\rule[.5ex]{2.5ex}{0.5pt}}

% \usepackage{algorithm}
% \usepackage{algpseudocode}

\newcommand{\mvar}[1]{\mathit{#1}}
\newcommand{\qt}[1]{\mbox{`#1'}}
% \algrenewcommand{\algorithmicreturn}{\State \textbf{return}}
\usepackage{tcolorbox}
\usepackage{subfigure} 

% For citations
\usepackage{natbib}

% For algorithms
\usepackage{algorithm}
\usepackage{algorithmic}

% disable the below lines if you want auto-indent
\newlength\tindent
\setlength{\tindent}{\parindent}
\setlength{\parindent}{0pt}
\renewcommand{\indent}{\hspace*{\tindent}}